% Pro Q. Ligario --- headnote
% pro-ligario, November 46 BC, Rome

\textit{Delivered before Julius Caesar in the Forum in late
November 46 BC, the second of Cicero's three so-called
``Caesarian'' speeches --- after \textit{Pro Marcello},
delivered in the Senate two months earlier, and before
\textit{Pro Rege Deiotaro}, delivered the following year in
Caesar's own house. The occasion was not a Senate motion
this time but a hearing in which Caesar sat as judge. Q.
Tubero, son of the jurist L. Tubero, had laid a charge
against Q. Ligarius, a Pompeian who had served as legate to
the proconsul C. Considius in Africa, who on Considius's
departure had been forced by the provincials to take the
province in hand, who had then refused the incoming Pompeian
commander P. Attius Varus little and so was held to have been
on the African Pompeian side that fell at Thapsus. Ligarius
was already in exile when Tubero, himself a former Pompeian
pardoned by Caesar after Pharsalus, brought the prosecution.
Cicero's task was to turn the case from a forensic charge
into a petition for clemency.}

\textit{The speech is more pointed and more openly forensic
than \textit{Pro Marcello}. Where the earlier speech is an
encomium of Caesar's \textit{clementia} addressed to a
Senate that has just granted a pardon, \textit{Pro Ligario}
must wring one out of a Caesar who has not yet decided. The
strategy is to confess the charge --- Ligarius was indeed in
Africa --- and then turn the prosecution back on Tubero. The
prosopopoeic attack of \textsection 9 is one of the most
famous passages in Cicero: \textit{quid enim, Tubero, tuus
ille destrictus in acie Pharsalica gladius agebat? cuius
latus ille mucro petebat?} --- ``what was that drawn sword
of yours doing in the line at Pharsalus? whose flank was
that point seeking?'' Tubero, who is himself a pardoned
Pompeian, has no answer that does not also acquit Ligarius.
\textsection 11 turns the screw further: the prosecution is
not for condemnation, since Ligarius is already exiled, but
for death --- and that, Cicero says, is the style of
``frivolous Greeks or savage barbarians,'' not of any Roman
citizen before this hour.}

\textit{The whole arc moves toward the deification maxim of
\textsection 38: \textit{homines enim ad deos nulla re
propius accedunt quam salutem hominibus dando} --- ``for men
come closer to the gods by nothing than by giving safety to
men.'' \textsection 30 contrasts the forensic register
(``gentlemen of the jury; he erred, he slipped'') with the
filial register Cicero claims for this occasion: he is not
pleading before a judge but before a parent. \textsection
32--34 broadens the appeal --- bringing forward the mourning
dress of T. Brocchus, the three Ligarii brothers, the wider
Sabine country --- and \textsection 33 turns one of
Caesar's own sayings against him: ``you considered all yours
who were not against you.'' The speech is short by Ciceronian
standards (about thirty-eight sections), but the architecture
is tight; ancient critics already singled it out as a
specimen of forensic art, and tradition has it that Caesar
read it through with mounting agitation and pardoned Ligarius
on the spot.}
